%% ===========================================================================
\section{Prompt of the LLM judge and questions asked to Laya}
\label{app:prompts}
%% ===========================================================================
%% Copied verbatim from laya_slots/judges/llm_judge.py (PROMPT_VERSION = "slot-judge-v1")
%% and laya_slots/judges/laya_judge.py. Do not edit by hand: re-copy if the code changes.

\subsection{LLM judge (prompt \texttt{slot-judge-v1})}

\paragraph{System message}
\begin{verbatim}
You are a slot-filling judge for a task-oriented dialogue system. You answer
with a single capital letter and nothing else.
\end{verbatim}

\paragraph{User message} Fields in braces are filled per request; \texttt{\{options\}} lists the candidates as ``\texttt{A. value}'', ``\texttt{B. value}''\ldots, followed by the ``not specified'' option. Empty fields are rendered as \texttt{unknown} (intent), \texttt{none} (filled slots) and \texttt{(none)} (history).
\begin{verbatim}
Task: choose the value that the user's last message gives for the slot below.
Pick the candidate that the user expresses or clearly refers to in the last
message (the history and the slots already filled are context only). If the
last message gives no value for this slot, or none of the candidates is right,
pick the "not specified" option.
The conversation may be in a language other than English; candidates are
written as they appear.

Slot: {slot}
Slot description: {description}
Current intent: {intent}
Slots already filled: {filled}
Conversation history:
{history}

User's last message ({lang}): {utterance}

Candidates:
{options}

Answer with the letter of one candidate only.
\end{verbatim}
(Line breaks inside paragraphs are added here for the page width; in the code each paragraph is one line.)

\paragraph{``Not specified'' option}
\begin{verbatim}
not specified (the last message gives no value for this slot)
\end{verbatim}

\paragraph{Few-shot format} (not used in the reported runs) Each shot is the same user message followed by an assistant turn containing the letter; in plain-completion mode the answer is rendered as \texttt{Answer: \{letter\}}.

\paragraph{Reading the answer} One token is generated at temperature 0; the top-20 log-probabilities of the first token are read; the probability masses of the tokens that reduce to an option letter (after removing the SentencePiece and byte-level space markers and trailing ``\texttt{.}'', ``\texttt{)}'', ``\texttt{:}'') are summed and renormalised over the offered letters. Reasoning is disabled with \texttt{enable\_thinking: false} for the Qwen3.8 and Gemma~4 chat templates.

\subsection{Laya questions}

\paragraph{Single mode (main condition)} Question of type \texttt{choice}; options are ``\texttt{A: <candidate>}'', ``\texttt{B: <candidate>}''\ldots, then the ``not specified'' option.
\begin{verbatim}
Which value does the user's last message give for the slot "{slot}"
({description})?
\end{verbatim}
``Not specified'' option:
\begin{verbatim}
not specified: the last message gives no value for this slot
\end{verbatim}

\paragraph{Two-step mode (implemented, not evaluated in this study)} First question (``is the slot expressed?''), of type \texttt{noul} (true/false) or, as a fallback against the bias of that question type, \texttt{choice} with neutral keys A/B:
\begin{verbatim}
Does the user's last message give a value for the slot "{slot}"
({description})?
  A / true:  yes, the last message gives a value for this slot
  B / false: no, the last message gives no value for this slot
\end{verbatim}
Second question, over the candidates only (no ``not specified'' option), in the same forward pass:
\begin{verbatim}
The user's last message gives a value for the slot "{slot}" ({description}).
Which one?
\end{verbatim}
The final distribution is $p(c_i)=P(\text{expressed})\cdot p(c_i\mid\text{expressed})$ and $p(\none)=1-P(\text{expressed})$; below a threshold on $P(\text{expressed})$ (0.5 by default, to be tuned on development data) the judge answers ``not specified''.

\paragraph{State} A dictionary with the fields \texttt{last\_user\_message}, \texttt{intent}, \texttt{slots\_already\_filled} (JSON, multi-turn only) and \texttt{history} (``\texttt{speaker: text}'' lines), in this order.
